% Edited conversion of source p. 38.
\ReportHeading
  {Дослідження впливу випадкових наскрізних пошкоджень поверхні на несучу здатність моделей циліндричних оболонок}
  {А.~П. Дзюба, Л.~Д. Левитіна, А.~Г. Пацюк}
  {Дніпровський національний університет імені Олеся Гончара}
  {}

Наведено результати експериментальних досліджень і чисельного скінченно-елементного аналізу втрати несучої здатності поздовжньо стиснутих моделей циліндричних оболонок із наскрізними пошкодженнями.

Концентратори напружень у вигляді тріщин, отворів і пробоїн різної форми створювали ударним навантаженням різною кількістю твердих фрагментів, варіюючи їх масу, форму, площу контакту, швидкість та кут проходження крізь поверхню моделі. Досліджено вплив цих параметрів на характер деформування та критичну руйнівну силу, за якої оболонка втрачала несучу здатність.

Розроблено методику системних експериментальних досліджень, яка дає змогу з мінімальними витратами варіювати параметри процесу і визначати найнебезпечніші поєднання зовнішнього навантаження та типу пошкодження. Використано створений авторами лабораторний комплекс, до складу якого входять пристрої статичного навантаження моделей на розтяг--стиск, згин і кручення, засоби ударного навантаження, а також обладнання для спостереження та вимірювання основних параметрів руйнування.

У випадках, де пряме зіставлення було можливим, отримано добрий збіг експериментальних результатів із чисельним скінченно-елементним аналізом. Результати створюють основу для методології прогнозування несучої здатності, або живучості, попередньо навантажених оболонкових конструкцій після виникнення наскрізних пошкоджень різної природи.

\textit{Дослідження виконано за підтримки Національного фонду досліджень України, проєкт №~2025.07/0117.}
